\scanpage{15}

on obtient des bijections inverses l'une de l'autre
\diagram{\begin{tikzcd}[ampersand replacement=\&]
\pi \arrow[r,bend left=15,"\exp"] \& \gamma \arrow[l,bend left=15,"\log"]
\end{tikzcd}}
Bien entendu, $\pi$ est un idéal \inserted{(nilpotent)} de $L_\Lambda(M)$, et
à fortiori une sous-algèbre de Lie. En
transportant à $\pi$ la structure de groupe de
$\gamma$, on trouve une structure de groupe \unsure{dans}
$\pi$, qui se déduit de sa structure
d'algèbre de Lie par la formule de
Hausdorff-Campbell. Si $M$ est une
$\Lambda$-Algèbre \inserted{$\unsure{B|A}$}, on trouve le

\uline{Corollaire} les \uline{automorphismes} de la
$\Lambda$-algèbre \uline{$\unsure{B|A}$} (où $\Lambda$ est une algèbre
$k$-augmentée : \deleted{\unsure{ayant}} idéal d'augmentation
nilpotent, $k$ est une $\mathbb Q$-algèbre)
qui induisent l'identité sur $B\otimes_\Lambda k$
correspondent biunivoquement aux
$\Lambda$-dérivations \inserted{$D$} de $B$ telles que $DB\subset I(\Lambda)B$.

\mspage{14}
